\section{Problem Definition}
\label{sec:problem}

Consider a set of $N\!\ge\!4$ fixed anchors with positions
$\{\avec_k\}_{k=1}^{N}\subset\mathbb{R}^3$ expressed in a common world
frame. A rover carries $M\in\{1,2\}$ tags on a rigid plate; for $M{=}2$
the two tag antennas are separated by a known baseline $L$. The rover
motion is planar, so its configuration is the body pose
$\xvec=(c_x,c_y,\psi)\in SE(2)$---planar position $\mathbf{c}$ and
heading $\psi$---which is extended to the full six-degree-of-freedom
pose using the height $z$ and the roll and pitch $(\phi,\theta)$ read
from the on-board IMU. The two tag positions follow from the body pose
as
\begin{equation}
  \pvec_{1,2} = \mathbf{c} \pm \tfrac{L}{2}\,[\cos\psi,\ \sin\psi]^{\!\top},
  \label{eq:rigid}
\end{equation}
for the front and rear tag respectively.

Each double-sided two-way ranging exchange between tag $i$ and anchor
$k$ produces a range estimate that departs from the true geometric
distance by several additive terms,
\begin{equation}
  \rho_{ik} = \lVert \pvec_i - \avec_k \rVert
            + \delta_k + \gamma_k(P_{ik}) + \Delta_{ik} + \varepsilon_{ik},
  \label{eq:measmodel}
\end{equation}
where $\delta_k$ is the per-anchor antenna-delay bias, $\gamma_k(P_{ik})$
a bias that depends on the received power $P_{ik}$, $\Delta_{ik}\ge 0$
the excess path length introduced under NLOS conditions, and
$\varepsilon_{ik}\sim\mathcal{N}(0,\sigma^2)$ zero-mean measurement
noise. The tags report at different instants, so the ranges arrive
asynchronously rather than as synchronized snapshots.

The problem is twofold. First, given the anchor positions and a stream
of asynchronous ranges~\eqref{eq:measmodel}, estimate the body
trajectory $\{\xvec_t\}$ online at the sensor rate ($\ge\!5$\,Hz) while
suppressing the systematic biases and NLOS excess. Second, when the
anchor positions are not externally available, estimate
$\{\avec_k\}$ from inter-anchor ranging alone. The remainder of the
paper addresses these two problems and evaluates the resulting accuracy
against an independent optical ground truth.
